\documentclass[11pt]{article}
\usepackage[a4paper,margin=18mm]{geometry}
\usepackage[no-math]{fontspec}
\setmainfont{Latin Modern Roman}
\usepackage{amsmath,amssymb,amsthm,xfrac,xspace,hyperref,graphicx,comment}
\excludecomment{DefTralics}

\providecommand{\bibliofont}{\small}
\newcommand{\ie}{i.e.,\xspace}
\newcommand{\eg}{e.g.,\xspace}
\newcommand{\etc}{etc.\xspace}
\newcommand{\cf}{cf.\xspace}
\newcommand{\resp}{resp.\xspace}
\newcommand{\smaller}{\small}
\newcommand{\Subsection}{\subsection}
\newcommand{\Subsubsection}{\subsubsection}
\newcommand{\skpt}{\smallskip}
\emergencystretch=3em
\numberwithin{equation}{section}\providecommand{\dpl}{\displaystyle}
\input{author-preamble.tex}
\begin{document}
\begin{center}{\Large Stable item 2038}\end{center}
\paragraph{Source and attribution.}
Daniel Álvarez-Gavela, Kiyoshi Igusa, \emph{A Legendrian Turaev torsion via generating families}, Journal de l’École polytechnique — Mathématiques 8 (2021), publisher version, DOI 10.5802/jep.141. \href{https://jep.centre-mersenne.org/articles/10.5802/jep.141/}{Publisher article}. Authors' text: \href{https://creativecommons.org/licenses/by/4.0/}{CC BY 4.0}.
\paragraph{Editorial problem boundary.}
Prove only Theorem 3.36 below: Legendrian-isotopy invariance of the printed Turaev torsion sets for Euler-type Legendrians and closed torsion pairs, with every original orientability, connectedness, local-system and generating-family assumption retained. The complete weak-Euler-structure, fibre-independence, balanced even-stabilisation and restriction core is retained. Mesh classification, link examples and non-naturality assertions are context only; the named homotopy-lifting and Cerf/Morse prerequisites remain explicit references.
\paragraph{Difficulty (editorial): H2.}
Classical Morse torsion and the cited homotopy-lifting theorem do not alone identify the torsion of an Euler-type generating family independently of its choices. The author constructs the required weak Euler structures from the Legendrian components and proves invariance under the relevant fibrewise Cerf changes, balanced even stabilisation and restriction. Tracking the compatible path classes and unit indeterminacy through those constructions is the nonroutine core used by the final isotopy comparison.
\paragraph{Editorial transcription note.}
Original author language, mathematics and ordering are preserved. Publisher front matter is replaced by this independent article wrapper. Bibliography is recovered from matching cached publisher HTML, including its TeX formulas. No new proof or mathematical repair is supplied. Surrounding original results are retained for conservative context and are not extra selected corpus claims.
\clearpage
\input{author-body.tex}
\begin{thebibliography}{99}
\bibitem{AGV85} V. I. Arnold, S. M. Gusein-Zade \& A. N. Varchenko - Singularities of differentiable maps. Volume 1 , Modern Birkhäuser Classics , Birkhäuser/Springer, New York, 2012
\bibitem{AK18} M. Abouzaid \& T. Kragh - “Simple homotopy equivalence of nearby Lagrangians” , Acta Math. 220 (2018) no. 2, p. 207-237
\bibitem{B64} D. Barden - On the structure and classification of differential manifolds , Ph. D. Thesis, Cambridge University, 1964
\bibitem{5W} K. Baur, E. Faber, S. Gratz, K. Serhiyenko \& G. Todorov - “Friezes satisfying higher $\mathrm{SL}_k$-determinants” , 2018, to appear in Algebra \& Number Theory
\bibitem{BL95} J.-M. Bismut \& J. Lott - “Flat vector bundles, direct images and higher real analytic torsion” , J. Amer. Math. Soc. 8 (1995) no. 2, p. 291-363
\bibitem{C70} J. Cerf - “La stratification naturelle des espaces de fonctions différentiables réelles et le théorème de la pseudo-isotopie” , Publ. Math. Inst. Hautes Études Sci. 39 (1970), p. 5-173
\bibitem{C84} M. Chaperon - “Une idée du type ‘géodésiques brisées’ pour les systèmes hamiltoniens” , C. R. Acad. Sci. Paris Sér. I Math. 298 (1984) no. 13, p. 293-296
\bibitem{C17} F. Charette - “Quantum Reidemeister torsion, open Gromov-Witten invariants and a spectral sequence of Oh” , Internat. Math. Res. Notices (2019) no. 8, p. 2483-2518
\bibitem{C96} Y. V. Chekanov - “Critical points of quasi-functions and generating families of Legendrian manifolds” , Funct. Anal. Appl. 30 (1996) no. 2, p. 118-128
\bibitem{C02} Y. V. Chekanov - “Differential algebra of Legendrian links” , Invent. Math. 150 (2002) no. 3, p. 441-483
\bibitem{CM18} R. Casals \& E. Murphy - “Differential algebra of cubic planar graphs” , Adv. Math. 338 (2018), p. 401-446
\bibitem{dR39} G. de Rham - “Sur les complexes avec automorphismes” , Comment. Math. Helv. 12 (1940), p. 191-211
\bibitem{DR11} G. Dimitroglou Rizell - “Knotted Legendrian surfaces with few Reeb chords” , Algebraic Geom. Topol. 11 (2011) no. 5, p. 2903-2936
\bibitem{DWW03} W. Dwyer, M. Weiss \& B. Williams - “A parametrized index theorem for the algebraic $K$-theory Euler class” , Acta Math. 190 (2003) no. 1, p. 1-104
\bibitem{EES07} T. Ekholm, J. Etnyre \& M. G. Sullivan - “Legendrian contact homology in $P\times \mathbb{R}$” , Trans. Amer. Math. Soc. 359 (2007) no. 7, p. 3301-3335
\bibitem{EG98} Y. M. Eliashberg \& M. Gromov - “Lagrangian intersection theory: finite-dimensional approach” , in Geometry of differential equations , Amer. Math. Soc. Transl. Ser. 2 , vol. 186 , American Mathematical Society, Providence, RI, 1998, p. 27-118
\bibitem{E98} Y. M. Eliashberg - “Invariants in contact topology” , in Proceedings of the ICM, Vol. II (Berlin, 1998) , Doc. Math. , Deutsche Mathematiker-Vereinigung, Berlin, 1998, p. 327-338, Extra Vol. II
\bibitem{EM12} Y. M. Eliashberg \& N. M. Mishachev - “The space of framed functions is contractible” , in Essays in mathematics and its applications , Springer, Heidelberg, 2012, p. 81-109
\bibitem{FI04} D. Fuchs \& T. Ishkhanov - “Invariants of Legendrian knots and decompositions of front diagrams” , Moscow Math. J. 4 (2004) no. 3, p. 707-717, 783
\bibitem{FR11} D. Fuchs \& D. Rutherford - “Generating families and Legendrian contact homology in the standard contact space” , J. Topology 4 (2011) no. 1, p. 190-226
\bibitem{F35} W. Franz - “Über die Torsion einer Überdeckung” , J. reine angew. Math. 173 (1935), p. 245-254
\bibitem{F95} K. Fukaya - “The symplectic $s$-cobordism conjecture: a summary” , in Geometry and physics (Aarhus, 1995) , Lecture Notes in Pure and Appl. Math. , vol. 184 , Dekker, New York, 1997, p. 209-219
\bibitem{GKS12} S. Guillermou, M. Kashiwara \& P. Schapira - “Sheaf quantization of Hamiltonian isotopies and applications to nondisplaceability problems” , Duke Math. J. 161 (2012) no. 2, p. 201-245
\bibitem{H04} M. B. Henry - “Connections between Floer-type invariants and Morse-type invariants of Legendrian knots” , Pacific J. Math. 249 (2011) no. 1, p. 77-133
\bibitem{HI19} E. J. Hanson \& K. Igusa - “A counterexample to the $\phi $-dimension conjecture” , 2019
\bibitem{HL99} M. Hutchings \& Y.-J. Lee - “Circle-valued Morse theory and Reidemeister torsion” , Geom. Topol. 3 (1999), p. 369-396
\bibitem{HR15} M. B. Henry \& D. Rutherford - “Equivalence classes of augmentations and Morse complex sequences of Legendrian knots” , Algebraic Geom. Topol. 15 (2015) no. 6, p. 3323-3353
\bibitem{HW73} A. Hatcher \& J. Wagoner - Pseudo-isotopies of compact manifolds , Astérisque , vol. 6 , Société Mathématique de France, Paris, 1973
\bibitem{K} K. Igusa - “The generalized Grassmann invariant” , preprint
\bibitem{I79} K. Igusa - The $\mathrm{Wh}_3(\pi )$ obstruction for pseudoisotopy , Ph. D. Thesis, Princeton University, 1979
\bibitem{I82} K. Igusa - “What happens to Hatcher and Wagoner’s formulas for $\pi _{0}C(M)$ when the first Postnikov invariant of $M$ is nontrivial?” , in Algebraic $K$-theory, number theory, geometry and analysis (Bielefeld, 1982) , Lect. Notes in Math. , vol. 1046 , Springer, Berlin, 1984, p. 104-172
\bibitem{I87} K. Igusa - “The space of framed functions” , Trans. Amer. Math. Soc. 301 (1987) no. 2, p. 431-477
\bibitem{I88} K. Igusa - “The stability theorem for smooth pseudoisotopies” , $K$-Theory 2 (1988) no. 1-2, p. 1-355
\bibitem{IK93a} K. Igusa - “The Borel regulator map on pictures. I. A dilogarithm formula” , $K$-Theory 7 (1993) no. 3, p. 201-224
\bibitem{I02} K. Igusa - Higher Franz-Reidemeister torsion , AMS/IP Studies in Advanced Math. , vol. 31 , American Mathematical Society, Providence, RI, 2002
\bibitem{I04} K. Igusa - “Combinatorial Miller-Morita-Mumford classes and Witten cycles” , Algebraic Geom. Topol. 4 (2004), p. 473-520
\bibitem{I05} K. Igusa - Higher complex torsion and the framing principle , Mem. Amer. Math. Soc. , vol. 177, no. 835 , American Mathematical Society, Providence, RI, 2005
\bibitem{IK93} K. Igusa \& J. Klein - “The Borel regulator map on pictures. II. An example from Morse theory” , $K$-Theory 7 (1993) no. 3, p. 225-267
\bibitem{J89} S. M. Jekel - “A simplicial formula and bound for the Euler class” , Israel J. Math. 66 (1989) no. 1-3, p. 247-259
\bibitem{JKS16} B. T. Jensen, A. D. King \& X. Su - “A categorification of Grassmannian cluster algebras” , Proc. London Math. Soc. (3) 113 (2016) no. 2, p. 185-212
\bibitem{T06} J. Jordan \& L. Traynor - “Generating family invariants for Legendrian links of unknots” , Algebraic Geom. Topol. 6 (2006), p. 895-933
\bibitem{K89} J. R. Klein - The cell complex construction and higher R-torsion for bundles with framed Morse functions , Ph. D. Thesis, Brandeis University, 1989
\bibitem{K92} M. Kontsevich - “Intersection theory on the moduli space of curves and the matrix Airy function” , Comm. Math. Phys. 147 (1992) no. 1, p. 1-23
\bibitem{K18} T. Kragh - “Generating families for Lagrangians in $\mathbb{R}^{2n}$ and the Hatcher-Waldhausen map” , 2018
\bibitem{La11} F. Laudenbach - Transversalité, courants et théorie de Morse , Éditions de l’École polytechnique, Palaiseau, 2012
\bibitem{L05a} Y.-J. Lee - “Reidemeister torsion in Floer-Novikov theory and counting pseudo-holomorphic tori. I” , J. Symplectic Geom. 3 (2005) no. 2, p. 221-311
\bibitem{L05b} Y.-J. Lee - “Reidemeister torsion in Floer-Novikov theory and counting pseudo-holomorphic tori. II” , J. Symplectic Geom. 3 (2005) no. 3, p. 385-480
\bibitem{LS85} F. Laudenbach \& J.-C. Sikorav - “Persistance d’intersection avec la section nulle au cours d’une isotopie hamiltonienne dans un fibré cotangent” , Invent. Math. 82 (1985) no. 2, p. 349-357
\bibitem{M63} B. Mazur - “Relative neighborhoods and the theorems of Smale” , Ann. of Math. (2) 77 (1963), p. 232-249
\bibitem{M61} J. Milnor - “Two complexes which are homeomorphic but combinatorially distinct” , Ann. of Math. (2) 74 (1961), p. 575-590
\bibitem{M66} J. Milnor - “Whitehead torsion” , Bull. Amer. Math. Soc. 72 (1966), p. 358-426
\bibitem{MT96} G. Meng \& C. H. Taubes - “$\underline{\mathit{SW}}=$ Milnor torsion” , Math. Res. Lett. 3 (1996) no. 5, p. 661-674
\bibitem{E12} E. Murphy - “Loose Legendrian embeddings in high dimensional contact manifolds” , 2019
\bibitem{R35} K. Reidemeister - “Homotopieringe und Linsenräume” , Abh. Math. Sem. Univ. Hamburg 11 (1935) no. 1, p. 102-109
\bibitem{RS71} D. B. Ray \& I. M. Singer - “$R$-torsion and the Laplacian on Riemannian manifolds” , Adv. Math. 7 (1971), p. 145-210
\bibitem{RS18} D. Rutherford \& M. G. Sullivan - “Generating families and augmentations for Legendrian surfaces” , Algebraic Geom. Topol. 18 (2018) no. 3, p. 1675-1731
\bibitem{S05} J. M. Sabloff - “Augmentations and rulings of Legendrian knots” , Internat. Math. Res. Notices (2005) no. 19, p. 1157-1180
\bibitem{S06} J. M. Sabloff - “Duality for Legendrian contact homology” , Geom. Topol. 10 (2006), p. 2351-2381
\bibitem{S86} J.-C. Sikorav - “Sur les immersions lagrangiennes dans un fibré cotangent admettant une phase génératrice globale” , C. R. Acad. Sci. Paris Sér. I Math. 302 (1986) no. 3, p. 119-122
\bibitem{S61} S. Smale - “Generalized Poincaré’s conjecture in dimensions greater than four” , Ann. of Math. (2) 74 (1961), p. 391-406
\bibitem{SS16} J. M. Sabloff \& M. G. Sullivan - “Families of Legendrian submanifolds via generating families” , Quantum Topol. 7 (2016) no. 4, p. 639-668
\bibitem{STWZ15} V. Shende, D. Treumann, H. Williams \& E. Zaslow - “Cluster varieties from Legendrian knots” , Duke Math. J. 168 (2019) no. 15, p. 2801-2871
\bibitem{STZ17} V. Shende, D. Treumann \& E. Zaslow - “Legendrian knots and constructible sheaves” , Invent. Math. 207 (2017) no. 3, p. 1031-1133
\bibitem{S02} M. G. Sullivan - “$K$-theoretic invariants for Floer homology” , Geom. Funct. Anal. 12 (2002) no. 4, p. 810-872
\bibitem{S17} L. S. Suárez - “Exact Lagrangian cobordism and pseudo-isotopy” , Internat. J. Math. 28 (2017) no. 8, p. 1750059, 35
\bibitem{T01} L. Traynor - “Generating function polynomials for Legendrian links” , Geom. Topol. 5 (2001), p. 719-760
\bibitem{T86} V. G. Turaev - “Reidemeister torsion in knot theory” , Uspehi Mat. Nauk 41 (1986) no. 1(247), p. 97-147, 240
\bibitem{T98} V. G. Turaev - “A combinatorial formulation for the Seiberg-Witten invariants of $3$-manifolds” , Math. Res. Lett. 5 (1998) no. 5, p. 583-598
\bibitem{V92} C. Viterbo - “Symplectic topology as the geometry of generating functions” , Math. Ann. 292 (1992) no. 4, p. 685-710
\bibitem{W76} J. B. Wagoner - “Diffeomorphisms, $K_{2}$, and analytic torsion” , Algebraic and geometric topology (Proc. Sympos. Pure Math., Stanford Univ., Stanford, Calif., 1976), Part 1 (Proc. Sympos. Pure Math.) XXXII (1978), p. 23-33
\bibitem{W82} F. Waldhausen - “Algebraic $K$-theory of spaces, a manifold approach” , in Current trends in algebraic topology, Part 1 (London, Ont., 1981) , CMS Conf. Proc. , vol. 2 , American Mathematical Society, Providence, RI, 1982, p. 141-184
\bibitem{W50} J. H. C. Whitehead - “Simple homotopy types” , Amer. J. Math. 72 (1950), p. 1-57
\end{thebibliography}
\end{document}
